% ------------------------------------------------------------------------
% Qualificação de Mestrado — PPG em Ciência da Computação — UFABC
% Documento principal: reúne o preâmbulo e todos os elementos do texto.
% Compilar com: latexmk -pdf main.tex   (ou pdflatex + bibtex + pdflatex x2)
% ------------------------------------------------------------------------
\documentclass[
	12pt,			% tamanho da fonte
	oneside,		% impressão em um só lado da folha (padrão para qualificação)
	a4paper,		% tamanho do papel
	chapter=TITLE,	% títulos de capítulos em CAIXA ALTA
	section=TITLE,	% títulos de seção em CAIXA ALTA
	brazil			% idioma principal do documento
	]{abntex2}

% ------------------------------------------------------------------------
% Preâmbulo: pacotes e configurações gerais do documento
% ------------------------------------------------------------------------

% ---
% Pacotes fundamentais
% ---
\usepackage{lmodern}			% Fonte Latin Modern
\usepackage[T1]{fontenc}		% Codificação de fonte
\usepackage[utf8]{inputenc}	% Codificação do arquivo fonte
\usepackage{lastpage}			% Número da última página
\usepackage{indentfirst}		% Indenta o primeiro parágrafo de cada seção
\usepackage{nomencl}			% Lista de abreviaturas e siglas
\usepackage{color}				% Suporte a cores
\usepackage{graphicx}			% Inclusão de figuras
\usepackage{microtype}			% Melhorias de justificação/hifenização
\usepackage{csquotes}			% Citações em bloco

% ---
% Pacotes de citação e referências (ABNT NBR 10520 e NBR 6023)
% ---
\usepackage[
	alf,				% estilo de citação autor-data (alfabético)
	abnt-etal-list=3,
	abnt-emphasize=bf	% negrito em vez de itálico para os títulos
	]{abntex2cite}

% ---
% Pacotes matemáticos e de tabelas (comuns em trabalhos de Ciência da Computação)
% ---
\usepackage{amsmath,amssymb}
\usepackage{booktabs}
\usepackage{multirow}
\usepackage{longtable}
\usepackage{algorithm}
\usepackage{algpseudocode}

% ---
% Hyperlinks: a classe abntex2 já carrega o pacote hyperref internamente
% (\RequirePackage{hyperref} em abntex2.cls), então aqui apenas ajustamos
% suas opções via \hypersetup, sem recarregar o pacote.
% ---
\usepackage[all]{hypcap}		% corrige links para figuras/tabelas

\hypersetup{
	pdftitle={\@title},
	pdfauthor={\@author},
	pdfcreator={LaTeX with abnTeX2},
	pdfkeywords={geração de moléculas}{SMILES}{recombinação}{ciência da computação},
	colorlinks=true,
	linkcolor=black,
	citecolor=black,
	filecolor=black,
	urlcolor=black,
	bookmarksdepth=4,
	pdfstartview=FitH
}

% ---
% Espaçamentos entre linhas e parágrafos (NBR 14724)
% ---
\SingleSpacing % espaçamento simples será sobrescrito por \OnehalfSpacing no textual
\setlength{\parindent}{1.3cm}
\setlength{\parskip}{0.2cm}

% ---
% Configuração dos títulos de capítulos e seções
% ---
\renewcommand{\ABNTEXchapterfontsize}{\Large}
\renewcommand{\ABNTEXsectionfontsize}{\large}
\renewcommand{\ABNTEXsubsectionfontsize}{\normalsize}

% ---
% Numeração de páginas no canto superior direito (padrão abntex2)
% ---
\makeatletter
\let\hrulefillorig\hrulefill
\makeatother


% ---
% Informações do documento (usadas na capa, folha de rosto e metadados do PDF)
% ---
\titulo{Geração combinatória de moléculas via recombinação de fragmentos representados em SMILES}
\autor{Crhisângela G. R. Ferreira}
\local{Santo André, SP}
\data{2026}
\instituicao{%
	Universidade Federal do ABC -- UFABC
	\par
	Programa de Pós-Graduação em Ciência da Computação}
\tipotrabalho{Qualificação de Mestrado}
% Preambulo com as informações do exemplar (usadas na folha de rosto)
\preambulo{Exame de qualificação apresentado ao Programa de
	Pós-Graduação em Ciência da Computação da Universidade Federal
	do ABC como parte dos requisitos para obtenção do título de
	Mestre em Ciência da Computação.}
\orientador[Orientador]{Prof. Dr. Jesús P. Mena-Chalco}
\coorientador[Coorientador]{Prof. Dr. Mauricio Coutinho}

\begin{document}

% ---------------------------------------------------------------
% Elementos pré-textuais
% ---------------------------------------------------------------
\pretextual

\imprimircapa
\imprimirfolhaderosto

% ============================================================
% RESUMO (NBR 6028) — parágrafo único, 150-500 palavras.
% Escrever por último, depois do restante do texto pronto.
%
% Tópicos a cobrir, nesta ordem:
%   - contextualização (1-2 frases)
%   - objetivo geral do trabalho
%   - metodologia empregada (resumida)
%   - resultados preliminares (já obtidos até a qualificação)
%   - próximos passos até a defesa
% ============================================================
\setlength{\absparsep}{18pt}
\begin{resumo}[Português]

% TODO: escrever o resumo.

\vspace{\onelineskip}

\noindent
\textbf{Palavras-chave}: palavra 1. palavra 2. palavra 3.
\end{resumo}

% ============================================================
% ABSTRACT — tradução do resumo em português. Escrever por último.
% ============================================================
\begin{resumo}[English]
 \begin{otherlanguage*}{english}

 % TODO: write the abstract (translation of "resumo.tex").

 \vspace{\onelineskip}

 \noindent
 \textbf{Keywords}: keyword 1. keyword 2. keyword 3.
 \end{otherlanguage*}
\end{resumo}


\pdfbookmark[0]{\listfigurename}{lof}
\listoffigures*
\cleardoublepage

\pdfbookmark[0]{\listtablename}{lot}
\listoftables*
\cleardoublepage

\begin{siglas}
	\item[ABNT] Associação Brasileira de Normas Técnicas
	\item[PPG] Programa de Pós-Graduação
	\item[SMILES] Simplified Molecular Input Line Entry System
	\item[QSAR] Quantitative Structure-Activity Relationship
	\item[UFABC] Universidade Federal do ABC
	\item[ZSA] Algoritmo de Zhang-Shasha (\emph{Zhang-Shasha Algorithm})
\end{siglas}


\pdfbookmark[0]{\contentsname}{toc}
\tableofcontents*
\cleardoublepage

% ---------------------------------------------------------------
% Elementos textuais
% ---------------------------------------------------------------
\textual

\chapter{Introdução}
\label{cap:introducao}

\section{Contextualização e motivação}
\label{sec:contextualizacao}

% Tópicos sugeridos:
%   - por que moléculas fotoativas/funcionais importam na prática
%     (ex.: OLEDs, células solares, singlet fission, corantes/fluoróforos)
%   - lacuna entre desenho racional de moléculas e exploração computacional
%     do espaço químico
%   - papel do SMILES como representação computacional de moléculas

\section{Definição do problema}
\label{sec:problema}

% Tópicos sugeridos:
%   - o que exatamente se quer gerar/otimizar (classe de moléculas-alvo)
%   - por que geração combinatória de fragmentos é o ponto de partida
%   - como diversidade/redundância estrutural entra como problema em aberto

\section{Organização do texto}
\label{sec:organizacao}

% TODO: um parágrafo curto remetendo a cada capítulo, escrito por último.

\chapter{Revisão de Literatura}
\label{cap:revisao-literatura}

\section{Fundamentos}
\label{sec:fundamentos}

% Tópicos sugeridos:
%   - representação SMILES (sintaxe, validação, canonização)
%   - noções de espectroscopia / propriedades ópticas de interesse
%   - TD-DFT e estados excitados (o que é, por que é usado para prever
%     propriedades fotoativas)
%   - representação de moléculas como grafos/árvores

\section{Estado da arte em geração molecular}
\label{sec:estado-da-arte}

% Tópicos sugeridos:
%   - algoritmos genéticos (GA) aplicados à geração/otimização molecular
%   - abordagens baseadas em grafos (GA-grafo, edição de grafos)
%   - modelos de linguagem sobre SMILES (RNN, transformers)
%   - autoencoders variacionais (VAE) para geração molecular
%   - busca em árvore / MCTS aplicada à construção molecular
%   - métricas de similaridade/diversidade estrutural (Tanimoto,
%     fingerprints, distância de edição de grafos/árvores, Zhang-Shasha)

\section{Síntese e lacunas}
\label{sec:sintese-lacunas}

% Tópicos sugeridos:
%   - quadro comparativo dos trabalhos discutidos
%     (abordagem x controle de geração x interpretabilidade x custo)
%   - lacuna que este trabalho pretende preencher

\chapter{Justificativa da Escolha Metodológica}
\label{cap:justificativa-metodologica}

% Tópicos sugeridos:
%   - por que uma abordagem baseada em árvores/grafos e recombinação
%     de fragmentos (GA-grafo), e não apenas um gerador de linguagem
%     (SMILES-RNN, VAE)
%   - vantagens: controlabilidade, interpretabilidade, determinismo,
%     rastreabilidade até os fragmentos de origem
%   - limitações reconhecidas dessa escolha frente aos modelos generativos
%     de aprendizado de máquina (ex.: cobertura do espaço químico)
%   - por que a distância de edição de árvores (Zhang-Shasha) é adequada
%     para medir diversidade estrutural nesse contexto

\chapter{Objetivos}
\label{cap:objetivos}

\section{Objetivo geral}
\label{sec:objetivo-geral}

% TODO: uma frase, verbo no infinitivo (ex.: "Desenvolver...", "Propor...").

\section{Objetivos específicos}
\label{sec:objetivos-especificos}

% Tópicos sugeridos (ajustar/expandir/reordenar):
\begin{enumerate}
	\item % geração combinatória de moléculas a partir de fragmentos
	\item % validação química / garantia de moléculas válidas
	\item % representação estrutural (árvore/grafo) das moléculas geradas
	\item % métrica de diversidade estrutural (Zhang-Shasha / GA-grafo)
	\item % comparação com abordagens da literatura (SMILES-RNN, VAE, GA)
\end{enumerate}

\chapter{Metodologia Proposta}
\label{cap:metodologia}

% Descrever o pipeline completo, mesmo que parcialmente implementado.

\section{Visão geral do pipeline}
\label{sec:visao-geral-pipeline}

% Tópicos sugeridos:
%   - diagrama/figura das etapas do pipeline
%   - o que já está implementado x o que é proposto/planejado

\section{Geração de fragmentos e combinação}
\label{sec:etapa-fragmentos}

% Tópicos sugeridos:
%   - categorias de fragmentos utilizadas (anéis, auxocromos, pontes, linhas)
%   - regras de combinação/recombinação
%   - validação química e deduplicação

\section{Representação estrutural (árvore/grafo)}
\label{sec:representacao-estrutural}

% Tópicos sugeridos:
%   - como a molécula (SMILES) é convertida em árvore/grafo
%   - papel do algoritmo de Zhang-Shasha / GA-grafo nessa representação

\section{Avaliação de diversidade estrutural}
\label{sec:avaliacao-diversidade}

% Tópicos sugeridos:
%   - métrica proposta (distância de edição de árvores)
%   - como será usada: caracterização do conjunto gerado, filtragem de
%     redundância, comparação com baseline

\section{Ambiente e ferramentas}
\label{sec:ambiente-ferramentas}

% Tópicos sugeridos:
%   - linguagem, bibliotecas (RDKit, pandas etc.), versões

\chapter{Resultados Preliminares}
\label{cap:resultados-preliminares}

% O gerador de SMILES já implementado entra aqui como prova de conceito
% de infraestrutura (não como resultado científico final).

\section{Infraestrutura de geração implementada}
\label{sec:resultados-infraestrutura}

% Tópicos sugeridos:
%   - o que já está funcionando (pipeline de geração combinatória)
%   - quantas moléculas válidas foram geradas
%   - taxa de aromatização / taxa de rejeição por invalidade química
%   - taxa de duplicatas removidas

\section{Diversidade estrutural observada}
\label{sec:resultados-diversidade}

% Tópicos sugeridos:
%   - resultados preliminares (se houver) da métrica de diversidade
%   - exemplos ilustrativos de moléculas geradas

\section{Limitações observadas e próximos passos}
\label{sec:resultados-limitacoes}

% Tópicos sugeridos:
%   - gargalos identificados (ex.: custo computacional)
%   - o que motiva as próximas etapas (ver Cronograma)

\chapter{Cronograma}
\label{cap:cronograma}

% Preencher com o que falta e quando (até a defesa). Ajustar linhas/colunas
% conforme as atividades reais do projeto.

\begin{table}[htb]
	\centering
	\caption{Cronograma de atividades previstas.}
	\label{tab:cronograma}
	\begin{tabular}{lcccccccc}
		\toprule
		\textbf{Atividade} & \textbf{T1} & \textbf{T2} & \textbf{T3} & \textbf{T4} & \textbf{T5} & \textbf{T6} & \textbf{T7} & \textbf{T8} \\
		\midrule
		% Revisão de literatura            &   &   &   &   &   &   &   &   \\
		% Consolidação da metodologia      &   &   &   &   &   &   &   &   \\
		% Implementação da métrica de diversidade &   &   &   &   &   &   &   &   \\
		% Experimentos / avaliação         &   &   &   &   &   &   &   &   \\
		% Escrita da dissertação           &   &   &   &   &   &   &   &   \\
		% Defesa                           &   &   &   &   &   &   &   &   \\
		\bottomrule
	\end{tabular}
	\fonte{Elaborado pelo autor.}
\end{table}


% ---------------------------------------------------------------
% Elementos pós-textuais
% ---------------------------------------------------------------
\postextual

\bibliography{referencias}

\end{document}
